\documentclass[journal,trans]{IEEEtran}
\usepackage{graphicx} % Required for inserting images
\usepackage[utf8]{inputenc}
\usepackage{graphicx}
\usepackage{float}
\usepackage{tikz}

\usepackage{cite}
\usepackage{amsmath,amssymb,amsfonts}
\usepackage{graphicx}
\usepackage{textcomp}
\usepackage{xcolor}
\usepackage{booktabs}

\usepackage{tabularx}

\usetikzlibrary{positioning, arrows.meta, shapes.geometric, calc, fit}
\usepackage{url}
 \usepackage[colorlinks=true,
            linkcolor=black,   % color de secciones
            citecolor=blue,    % <-- color de las citas clickeables
            urlcolor=blue]{hyperref}
\usetikzlibrary{positioning}
\ifCLASSINFOpdf
\else
\fi
\hyphenation{op-tical net-works semi-conduc-tor}


\newcommand{\titlepaper}{Interconnect Modeling for
Performance Analysis of HW/SW Coherence Protocols}

\title{\titlepaper}

\author{
    \begin{minipage}[t]{0.3\textwidth}
        \centering
        \IEEEauthorblockN{Bolaños Lopez Randall Bryan}
        \IEEEauthorblockA{email: \url{yitanr@estudiantec.cr} \\
        Escuela de Ingeniería en Computadores\\
        Instituto Tecnológico de Costa Rica}
    \end{minipage}
    \hfill
    \begin{minipage}[t]{0.3\textwidth}
        \centering
        \IEEEauthorblockN{Molina Redondo Andres}
        \IEEEauthorblockA{email: \url{andresmolina@estudiantec.cr} \\
        Escuela de Ingeniería en Computadores\\
        Instituto Tecnológico de Costa Rica}
    \end{minipage}
    \hfill
    \begin{minipage}[t]{0.3\textwidth}
        \centering
        \IEEEauthorblockN{Rodríguez Cordero Christopher David}\\
        \IEEEauthorblockA{email: \url{chrodriguez@estudiantec.cr} \\
        Escuela de Ingeniería en Computadores\\
        Instituto Tecnológico de Costa Rica}
    \end{minipage}
}

\markboth{Bolaños, Molina, Rodríguez, \titlepaper}%
{Shell \MakeLowercase{\textit{et al.}}: Sistema TDM con DTMF}

\begin{document}

\maketitle
\begin{abstract}

This paper presents the design and implementation of a SystemC-based simulation
model of a cache-coherent interconnect for a shared-memory multiprocessor system
composed of four processing elements (PEs), each with a private L1 cache.
The simulator integrates a software layer---parallel workload memory traces---with
a hardware layer comprising a formal bus transaction model, parametric L1 cache
logic, and two cache coherence protocols: MSI (Write-Invalidate family) and
Firefly (Write-Update family). A formal set of five bus transactions is defined:
\texttt{BusRead}, \texttt{BusReadX}, \texttt{BusUpgrade}, \texttt{BusUpdate},
and \texttt{BusInvalidate}. Two parallel workloads are evaluated---Blocked Matrix
Multiplication under high contention and a Cryptography/Network Pipeline under
a producer-consumer pattern---plus an exploratory migratory access scenario.
Results show that under high contention MSI generates 39 invalidations with zero
updates, while Firefly replaces them with 114 broadcast updates at equal total
transaction count. In the pipeline scenario, Firefly reduces total bus
transactions by 22\% (18 to 14) by proactively updating remote caches and
avoiding re-fetches. The migratory pattern shows a 41\% transaction reduction
under Firefly. The simulation exports VCD waveforms, per-cache statistics, and
CSV data for reproducible analysis.

\end{abstract}

\begin{IEEEkeywords}
cache coherence, MSI protocol, Firefly protocol, SystemC, interconnect modeling,
shared memory, multiprocessor, bus transactions, Write-Invalidate, Write-Update
 
\end{IEEEkeywords}

\IEEEpeerreviewmaketitle

\section{Introducción}



Modern multiprocessor systems rely on shared-memory architectures in which
multiple processing cores communicate through a common interconnect. Ensuring
that all processors observe a consistent view of shared data is the fundamental
challenge addressed by cache coherence protocols \cite{nagarajan2020primer}.
The choice of protocol has a direct impact on bus traffic, memory bandwidth
consumption, and overall system performance.
 
Two dominant protocol families exist. Write-Invalidate protocols, such as
MSI \cite{archibald1986cache}, respond to a write by invalidating all remote
copies of the modified block, forcing subsequent readers to re-fetch it.
Write-Update protocols, such as Firefly \cite{thacker1988firefly}, instead
broadcast the new value to all sharers, keeping their copies valid at the
cost of higher per-write bus traffic.
 
This project implements a cycle-accurate SystemC \cite{ieee2024systemc}
simulation of a four-PE shared-memory system with private L1 caches and a
cache-coherent bus. The simulator integrates memory traces generated by parallel
workloads with the hardware model, enabling a direct comparison of MSI and
Firefly under realistic access patterns. A key contribution of this work is the
formalization of a typed bus transaction set that makes the distinction between
\texttt{BusReadX} (write miss from I) and \texttt{BusUpgrade} (write hit from S)
explicit, improving both correctness and observability of the simulation.
 
The remainder of this paper is organized as follows. Section~\ref{sec:theory}
provides the theoretical background. Section~\ref{sec:system} describes the
system architecture and design decisions. Section~\ref{sec:results} presents
and analyzes the experimental results. Section~\ref{sec:conclusions} concludes.

\section{Theoretical Background}
\label{sec:theory}
 
\subsection{Cache Coherence}
 
In a shared-memory multiprocessor, each processor has one or more levels of
private cache. When two processors hold copies of the same memory block and one
modifies it, the other's copy becomes stale. Cache coherence protocols enforce
the invariant that all processors observe a consistent value for every memory
location \cite{hennessy2017computer, nagarajan2020primer}. A coherence protocol
defines stable states for each cache block and transitions triggered by local
processor requests or bus snooping events \cite{solihin2016fundamentals}.
 
\subsection{MSI Protocol}
 
MSI is the baseline Write-Invalidate protocol \cite{archibald1986cache}.
Each cache block resides in one of three stable states:
 
\begin{itemize}
  \item \textbf{Modified (M):} Present only in this cache; written locally.
        The cache holds exclusive ownership.
  \item \textbf{Shared (S):} Clean copy; may be present in multiple caches.
  \item \textbf{Invalid (I):} Not usable; must be fetched before use.
\end{itemize}
 
A write from state I issues \texttt{BusReadX}, fetching the block and
invalidating all remote copies. A write from state S issues \texttt{BusUpgrade},
which invalidates remote copies without re-transferring data---an optimization
that reduces unnecessary bus traffic \cite{hennessy2017computer}.
 
\subsection{Firefly Protocol}
 
Firefly \cite{thacker1988firefly} is a Write-Update protocol. On a write, the
interconnect broadcasts the new value to all caches that hold the block via a
\texttt{BusUpdate} transaction. All sharers remain in state S. If no sharer
exists, the writing cache gains exclusive ownership (state M) without
broadcasting. This approach eliminates future read misses for existing sharers
at the cost of additional traffic per write \cite{culler1999parallel}.
 
\subsection{Bus Transaction Formalization}
 
A bus transaction is a typed message emitted by a cache to the interconnect,
specifying both the requested operation and the required coherence action.
Table~\ref{tab:transactions} defines the five transactions implemented in this
work.


\begin{table}[htbp]
\caption{Formal Bus Transaction Set}
\label{tab:transactions}
\renewcommand{\arraystretch}{1.2}
\begin{tabularx}{\columnwidth}{lX}
\toprule
\textbf{Transaction} & \textbf{Description and State Transition} \\
\midrule
\texttt{BusRead}
  & Read miss (I$\to$S). Fetches block from memory or a downgraded M-owner. \\
\texttt{BusReadX}
  & Write miss from I (I$\to$M). Fetches block; invalidates all remote copies. \\
\texttt{BusUpgrade}
  & Write hit from S (S$\to$M). Invalidates remote copies; no data transfer. \\
\texttt{BusUpdate}
  & Firefly write. Broadcasts new value to all sharers; they remain in S. \\
\texttt{BusInvalidate}
  & Internal signal emitted by the interconnect to force I on remote caches. \\
\bottomrule
\end{tabularx}
\end{table}
 
\subsection{False Sharing}
 
False sharing occurs when processors access distinct data items that share a
cache block \cite{hennessy2017computer}. Under Write-Invalidate, a write to any
byte invalidates the entire block in all other caches, even if other processors
accessed different bytes. This is the primary source of coherence overhead in
the matrix multiplication workload.
 
% ─────────────────────────────────────────────
\section{System Design}
\label{sec:system}
 
\subsection{Overall Architecture}
 
The simulated system consists of four PEs connected to a shared cache-coherent
bus, each with a private L1 cache, backed by shared main memory
\cite{garcia2026ce4302}. The entire hardware model is implemented in
SystemC/C++17 \cite{ieee2024systemc}. Python is used only for post-processing
and plot generation.
 
\subsection{Processing Elements}
 
Each PE is an \texttt{SC\_THREAD} that reads a memory trace file and issues
\texttt{READ} or \texttt{WRITE} requests to its local L1 cache. Trace lines
follow the format \texttt{<pe\_id> <READ|WRITE> <hex\_address>}. Each PE
filters only lines assigned to its identifier, enabling all four PEs to execute
in parallel from the same file.
 
\subsection{L1 Cache Design}
 
The L1 cache is a parametric set-associative module supporting LRU and FIFO
replacement. Table~\ref{tab:l1config} shows the default configuration used
for all experiments. This configuration was chosen to ensure sufficient
evictions and conflict misses to stress the coherence protocol
\cite{hennessy2017computer, patterson2013computer}.
 \begin{table}[htbp]
 \centering 
\caption{L1 Cache Default Configuration}
\label{tab:l1config}
\renewcommand{\arraystretch}{1.2}
\begin{tabular}{ll}
\toprule
\textbf{Parameter}          & \textbf{Value} \\
\midrule
Total size                  & 1\,024 bytes   \\
Block size                  & 64 bytes       \\
Associativity               & 2-way          \\
Replacement policy          & LRU            \\
L1 hit latency              & 1\,ns          \\
L1 miss detect latency      & 1\,ns          \\
\bottomrule
\end{tabular}
\end{table}
 
The \texttt{access(op, addr)} method determines whether the block is present
and valid, tracks hit/miss/eviction/conflict-miss counters, and delegates to
the interconnect when a bus transaction is required.
 
\subsection{Interconnect and Transaction Model}
 
The interconnect exposes \texttt{issue\_transaction(tx, id, addr)} as its
primary interface. The cache selects the transaction type based on the current
block state and the active protocol:
 
\begin{itemize}
  \item Read miss $\Rightarrow$ \texttt{BusRead}
  \item Write from I, MSI $\Rightarrow$ \texttt{BusReadX}
  \item Write from S, MSI $\Rightarrow$ \texttt{BusUpgrade}
  \item Any write, Firefly $\Rightarrow$ \texttt{BusUpdate}
\end{itemize}
 
Internally, each transaction type is handled by a dedicated method
(\texttt{do\_bus\_read}, \texttt{do\_bus\_read\_x}, etc.), separating concerns
and making the protocol logic auditable. The interconnect timing model uses four
configurable parameters: bus arbitration (1\,ns), memory latency (5\,ns),
shared cache latency (2\,ns), and coherence overhead (2\,ns).

% Arbitraje FCFS
\paragraph{Arbitraje del Bus (FCFS)}
El interconect incluye un arbitro interno que implementa una política
first-come, first-served (FCFS). Las caches encolan sus solicitudes en una
estructura de peticiones dentro de `Interconnect` y un `SC_THREAD` de
concesión procesa la cola, otorgando el bus a una única petición por vez.
Esta implementación proporciona un ordenamiento estricto FCFS y facilita el
análisis de fairness y latencias de espera.
 
\subsection{Workloads}
 
\textbf{Workload 1 -- Blocked Matrix Multiplication:}
Four PEs compute sub-blocks in parallel and write to a shared accumulator
address. All PEs write to the same cache block, maximizing false sharing and
coherence contention.
 
\textbf{Workload 2 -- Cryptography/Network Pipeline:}
A four-stage pipeline where PE0 writes a packet, PE1 reads and rewrites,
PE2 reads, and the pattern repeats for multiple packets. This models a
unidirectional producer-consumer migration.
 
\textbf{Exploratory -- Migratory Pattern:}
A single shared token address is read and written by each PE in round-robin,
representing the worst case for Write-Invalidate protocols.
 
% ─────────────────────────────────────────────
\section{Results and Analysis}
\label{sec:results}
 
\subsection{Workload 1: Blocked Matrix Multiplication}
 
Table~\ref{tab:matmul} summarizes bus statistics for this high-contention
workload.
 \begin{table}[htbp]
 \centering
\caption{Bus Statistics -- Blocked Matrix Multiplication}
\label{tab:matmul}
\renewcommand{\arraystretch}{1.2}
\begin{tabular}{lrr}
\toprule
\textbf{Metric}         & \textbf{MSI} & \textbf{Firefly} \\
\midrule
Total transactions      & 44  & 44  \\
BusRead                 & 4   & 4   \\
BusReadX                & 40  & 0   \\
BusUpgrade              & 0   & 0   \\
BusUpdate               & 0   & 40  \\
Invalidations           & 39  & 0   \\
Broadcast updates       & 0   & 114 \\
Downgrades M$\to$S      & 0   & 0   \\
Simulation time (ns)    & 91  & 91  \\
\bottomrule
\end{tabular}
\end{table}
 
Under MSI, every write to the shared address issues \texttt{BusReadX}, producing
39 invalidations. Each PE must then re-fetch the block before its next write,
creating a ping-pong ownership transfer between all four caches---a classic
invalidation storm \cite{archibald1986cache}. Under Firefly, all 40 writes use
\texttt{BusUpdate}, producing 114 individual remote cache updates (up to 3 per
write). No invalidations occur and no re-fetches are needed. The total
transaction count is equal (44), but the coherence mechanism is entirely
different. Fig.~\ref{fig:invalidaciones} and Fig.~\ref{fig:updates} illustrate
this complementary behavior.
 
\subsection{Workload 2: Cryptography/Network Pipeline}
 
Table~\ref{tab:pipeline} shows results for the producer-consumer workload.
 \begin{table}[htbp]
 \centering
\caption{Bus Statistics -- Cryptography/Network Pipeline}
\label{tab:pipeline}
\renewcommand{\arraystretch}{1.2}
\begin{tabular}{lrr}
\toprule
\textbf{Metric}         & \textbf{MSI} & \textbf{Firefly} \\
\midrule
Total transactions      & 18  & 14  \\
BusRead                 & 8   & 4   \\
BusReadX                & 6   & 0   \\
BusUpgrade              & 4   & 0   \\
BusUpdate               & 0   & 10  \\
Invalidations           & 12  & 0   \\
Broadcast updates       & 0   & 17  \\
Downgrades M$\to$S      & 6   & 1   \\
Simulation time (ns)    & 46  & 46  \\
\bottomrule
\end{tabular}
\end{table}
 
A notable result is the appearance of 4 \texttt{BusUpgrade} transactions in
MSI. PE1 reads a block (entering state S) before writing it in the next stage.
Since it already holds valid data, only an upgrade is issued---no data
re-transfer occurs. This confirms the importance of distinguishing
S$\to$M upgrades from I$\to$M read-exclusive requests in the transaction model.
 
Firefly reduces total transactions from 18 to 14 (22\%). Proactive updates allow
downstream consumers to find data already valid in their caches, eliminating
four \texttt{BusRead} transactions that MSI requires after invalidation. This
is consistent with the theoretical advantage of Write-Update in producer-consumer
patterns \cite{nagarajan2020primer, thacker1988firefly}.
 
\subsection{Exploratory: Migratory Access Pattern}
 
Table~\ref{tab:migratorio} shows the migratory scenario results.
 \begin{table}[htbp]
 \centering
\caption{Bus Statistics -- Migratory Access Pattern}
\label{tab:migratorio}
\renewcommand{\arraystretch}{1.2}
\begin{tabular}{lrr}
\toprule
\textbf{Metric}         & \textbf{MSI} & \textbf{Firefly} \\
\midrule
Total transactions      & 61  & 36  \\
Invalidations           & 55  & 0   \\
Broadcast updates       & 0   & 96  \\
Downgrades M$\to$S      & 16  & 0   \\
Simulation time (ns)    & 91  & 73  \\
\bottomrule
\end{tabular}
\end{table}
 
This pattern is the worst case for MSI: every write invalidates all others,
forcing re-fetches before the next read. MSI generates 61 transactions vs.\
36 for Firefly---a 41\% reduction---and finishes 18\,ns sooner, consistent
with classical analysis of migratory patterns \cite{archibald1986cache,
culler1999parallel}. Fig.~\ref{fig:total_tx} and Fig.~\ref{fig:bandwidth}
visualize these differences across all three workloads.
 
\begin{figure}[htbp]
\centering
\includegraphics[width=\columnwidth]{figures/transacciones_totales.png}
\caption{Total bus transactions per workload and protocol.}
\label{fig:total_tx}
\end{figure}
 
\begin{figure}[htbp]
\centering
\includegraphics[width=\columnwidth]{figures/bus_tx_por_ns.png}
\caption{Approximate bus bandwidth (tx/ns) per workload and protocol.}
\label{fig:bandwidth}
\end{figure}
 
\begin{figure}[htbp]
\centering
\includegraphics[width=\columnwidth]{figures/invalidaciones.png}
\caption{Invalidations: MSI vs.\ Firefly (Firefly always zero).}
\label{fig:invalidaciones}
\end{figure}
 
\begin{figure}[htbp]
\centering
\includegraphics[width=\columnwidth]{figures/updates.png}
\caption{Broadcast updates: MSI vs.\ Firefly (MSI always zero).}
\label{fig:updates}
\end{figure}
 
% ─────────────────────────────────────────────
\section{Conclusions}
\label{sec:conclusions}
 
This work demonstrates that cache coherence protocol selection has a significant,
workload-dependent impact on bus traffic and execution time. Key findings are:
 
\begin{enumerate}
 
\item \textbf{MSI generates invalidation storms under high contention.}
In the matrix workload, 89\% of bus transactions (39/44) are invalidations that
force re-fetches, confirming the classical weakness of Write-Invalidate under
false sharing \cite{archibald1986cache}.
 
\item \textbf{Firefly trades invalidations for broadcast updates.}
Under high contention, Firefly eliminates all invalidations and replaces them
with 114 updates. Whether this is beneficial depends on the subsequent read
rate from the updated caches.
 
\item \textbf{Formalizing BusUpgrade vs.\ BusReadX improves correctness.}
The pipeline workload shows 4 \texttt{BusUpgrade} transactions that would have
been incorrectly counted as \texttt{BusReadX} without explicit state-based
transaction selection, avoiding spurious data transfers.
 
\item \textbf{Firefly outperforms MSI in producer-consumer and migratory
patterns.} Transaction reductions of 22\% and 41\% respectively confirm the
theoretical predictions of \cite{nagarajan2020primer, thacker1988firefly}.
 
\item \textbf{No protocol dominates all workloads.} Protocol selection must
be informed by the expected sharing pattern: Write-Update favors
producer-consumer and migratory access; Write-Invalidate may be preferable
for predominantly private or write-once workloads.
 
\end{enumerate}
 

\bibliographystyle{IEEEtran}
\begin{thebibliography}{00}
 
\bibitem{hennessy2017computer}
J.~L. Hennessy and D.~A. Patterson,
\emph{Computer Architecture: A Quantitative Approach}, 6th~ed.
Elsevier, 2017.
 
\bibitem{nagarajan2020primer}
V.~Nagarajan, D.~J. Sorin, M.~D. Hill, and D.~A. Wood,
\emph{A Primer on Memory Consistency and Cache Coherence}, 2nd~ed.
Morgan \& Claypool, 2020.
 
\bibitem{archibald1986cache}
J.~Archibald and J.-L. Baer,
``Cache coherence protocols: Evaluation using a multiprocessor simulation
model,''
\emph{ACM Trans. Comput. Syst.}, vol.~4, no.~4, pp.~273--298, Nov. 1986.
 
\bibitem{thacker1988firefly}
C.~P. Thacker, L.~C. Stewart, and E.~H. Satterthwaite~Jr.,
``Firefly: A multiprocessor workstation,''
\emph{IEEE Trans. Comput.}, vol.~37, no.~8, pp.~909--920, Aug. 1988.
 
\bibitem{ieee2024systemc}
IEEE,
\emph{IEEE Standard for Standard SystemC Language Reference Manual},
IEEE Std 1666-2023, 2024.
 
\bibitem{ieee2018systemverilog}
IEEE,
\emph{IEEE Standard for SystemVerilog --- Unified Hardware Design,
Specification, and Verification Language},
IEEE Std 1800-2017, 2018.
 
\bibitem{garcia2026ce4302}
R.~Garc\'ia~Fern\'andez and L.~Chavarr\'ia~Zamora,
``CE 4302 Proyecto Grupal 01: Modelado de Interconnect para el an\'alisis
de desempe\~no de protocolos de coherencia HW/SW,''
Instituto Tecnol\'ogico de Costa Rica, Cartago, Costa Rica, Mar. 2026.
 
\bibitem{solihin2016fundamentals}
Y.~Solihin,
\emph{Fundamentals of Parallel Multicore Architecture}.
CRC Press, 2016.
 
\bibitem{culler1999parallel}
D.~Culler, J.~P. Singh, and A.~Gupta,
\emph{Parallel Computer Architecture: A Hardware/Software Approach}.
Morgan Kaufmann, 1999.
 
\bibitem{patterson2013computer}
D.~A. Patterson and J.~L. Hennessy,
\emph{Computer Organization and Design: The Hardware/Software Interface},
5th~ed. Elsevier, 2013.
\end{thebibliography}

\end{document}